%%%%%%%%%%%%%%%%%%%%%%%%%%%%%%%%%%%%%%%%%%%%%%%%%%%%%%%%%%%%%%%%%%%
\section{The Bandwidth Tax, and Experiments to Measure It}
\label{sec:BandwidthTax}
%%%%%%%%%%%%%%%%%%%%%%%%%%%%%%%%%%%%%%%%%%%%%%%%%%%%%%%%%%%%%%%%%%%
% New Aug 28 pm, per MN's chat request: game out what a 10-20%
% bandwidth loss does to inference on an H100, and propose the tests.
% Parked as an appendix "for now" -- MN will decide whether it stays,
% moves into Sec. VII, or gets cut.

\inote[aiviolet]{parked here per your request (``fold this into a new
appendix for now about things we might want to try''); it can move
into Sec.~VII or get cut once you have read it}

What would it do to inference if monitoring consumed 10--20\% of an
H100's memory bandwidth? To first order the tax passes through
one-for-one, because decode is memory-bandwidth-bound across
essentially all practical serving configurations: throughput falls by
the stolen share, and inter-token latency rises as $1/(1-x)$, i.e.,
$+11\%$ at a 10\% tax and $+25\%$ at 20\%. In numbers, a 70B model at
fp8 is roughly \qty{70}{\giga\byte} of weights, so batch-1 decode on
one H100 has a streaming floor near 48~tokens/s -- about 43 under a
10\% tax and 38 under 20\%. At batch~32 with 8K contexts the weights
amortize but each stream drags its own key--value (KV) cache (roughly
\qty{1.3}{\giga\byte} per stream for a grouped-query 70B at fp8), so a
decode step moves $\sim$\qty{112}{\giga\byte} and the node serves
$\sim$950 aggregate tokens/s -- again falling in proportion, to
$\sim$760 at a 20\% tax. Notably, batching does not remove the tax,
since KV traffic keeps decode bandwidth-bound at every practical batch;
prefill is the exception, being compute-bound, so time-to-first-token
moves little.\footnote{Arithmetic on public specifications, not
measurements. Achievable HBM bandwidth is typically 85--90\% of the
\qty{3.35}{\tera\byte\per\second} figure, which shifts the absolute
numbers but not the proportional argument.}

Three second-order effects could make the same tax cost more than its
face value. {\bf (1)} \textbf{Queueing amplification:} a serving node
runs at some utilization against a latency target; removing 20\% of
capacity at a fixed arrival rate moves utilization of 0.75 to roughly
0.94, at which point tail latency grows several-fold and
goodput-under-SLO falls much faster than the bandwidth did. The unit
an operator would reason in is therefore fleet size: a 10--20\%
bandwidth tax is 10--25\% more GPUs to hold the same service level.
{\bf (2)} \textbf{Contention is not a clean partition:} the monitor's
bank scan is a streaming read that also evicts the \qty{50}{\mega\byte}
L2 that attention and weight kernels reuse, and DRAM scheduler and
row-buffer interference between co-running streams can push observed
degradation past the raw bandwidth share. Current parts offer no
fine-grained bandwidth QoS; MIG isolates the monitor, but a MIG slice
is itself a dedicated capacity loss. {\bf (3)} \textbf{Energy:} the
tax is approximately proportional in HBM access energy, which at
realistic $K$, layer counts, and token rates is the kilowatt-scale
arithmetic of the main paper's Sec.~VII-E.

Each of these maps to a short measurement, none requiring training:
{\bf (1)} a \emph{bandwidth-thief} microbenchmark -- co-run a tuned
streaming kernel (or the monitoring GEMV itself, over a $K\times d$
bank) against a serving stack (e.g., vLLM) running a 70B on one H100;
sweep the thief at 0/5/10/20/40\% of peak bandwidth; record decode
tokens/s, inter-token latency at the median and tail, and
time-to-first-token, at batch 1/8/32/64 and 1K/8K/32K contexts. The
slope of throughput against stolen bandwidth is the pass-through
factor: 1.0 confirms the first-order story, larger measures the
contention overshoot.
{\bf (2)} the same thief with L2 access-policy controls (persistence
windows, bypass hints), separating cache pollution from bandwidth.
{\bf (3)} the thief under MPS co-residence against a MIG slice,
pricing the isolation options.
{\bf (4)} the monitor end-to-end: a per-token FAISS/cuBLAS scan at
$K=10^{4}$--$10^{5}$, $d=8{,}192$ hooked into the serving loop, with
the measured tokens/s delta compared against the analytic
$K\,d\,b\,R$ of Eq.~(2) of the main paper.
{\bf (5)} a fixed-arrival-rate load test measuring goodput at an
inter-token-latency SLO with and without a 10/20\% thief, which is
the queueing-amplification number.
Experiment (1) is the highest value per day and pairs naturally with
the tuned-GPU-baseline item of the main paper's Sec.~VII-F;
a Jetson-class repeat of (1) at batch 1 covers the edge case, where
the first-order story should hold almost exactly.
